% ==========================================
% APPENDIX: NILPOTENT COMPLETION AND KOSZUL DUALITY
% For Non-Quadratic Chiral Algebras
% Status: FRONTIER appendix (Stratum II).
% The theorematic core of the monograph is the finite-type
% Koszul-locus theory of Chapters 2, 8, and 9.  This appendix
% records the leading extension of that core, not a replacement.
% ==========================================

\chapter{Nilpotent-completion frontier for non-quadratic Koszul duality}
\label{app:nilpotent-completion}


\begin{abstract}
This appendix develops a frontier completion package extending the
finite-type geometric Koszul formalism to a larger class of filtered
chiral algebras under finite-generation and polynomial-growth
hypotheses.  For non-quadratic chiral algebras such as Virasoro and
W-algebras, the bar construction produces an infinite coalgebra that
requires completion to yield a usable completed curved bar model and
dual shadow data.  The appendix records the $I$-adic framework, the
finite-stage infrastructure, and the exact missing lemmas that separate
the standard completed examples from a universal completion theorem.
\end{abstract}

\begin{remark}[Dependency warning]
\label{rem:nilpotent-completion-dependency}
The theorematic core of the monograph remains the finite-type
Koszul-locus theory of Chapters~\ref{chap:koszul-pairs}
and~\ref{chap:higher-genus}.  The present appendix records the leading
extension of that core, not a replacement of it.  Every result below
should be read as belonging to the completion frontier (Stratum~II),
and its validity is conditional on the finite-generation and
polynomial-growth hypotheses stated in
Theorem~\ref{thm:completion-convergence}.
\end{remark}

\section{The problem: why completion is necessary}

\subsection{Quadratic vs. non-quadratic}

\begin{observation}[Classical Koszul theory limitation]\label{obs:quadratic-limitation}
Classical Koszul duality applies to quadratic algebras—those presented by generators and quadratic relations. Examples:
\begin{itemize}
\item Symmetric algebra: Sym(V) with relations $xy = yx$
\item Exterior algebra: $\Lambda(V)$ with relations $xy = -yx$, $x^2 = 0$
\item Polynomial ring: $k[x_1, \ldots, x_n]$
\end{itemize}

However, many chiral algebras require completion beyond the quadratic
setting: the Virasoro OPE $T(z)T(w) \sim c/(z-w)^4 + \cdots$ is
\emph{inhomogeneous} quadratic (bilinear plus central term), requiring
curved Koszul duality (\`{a}~la Positselski) and, in the manuscript's
main line, producing only the same-family shadow
$\mathrm{Vir}_{26-c}$ rather than an H-level infinite-generator dual
object; W-algebras ($N \geq 3$) involve cubic and higher-order OPE
relations; and the affine Yangian has relations involving spectral
parameters.

For these, the bar construction $\bar{B}(\cA)$ remains a perfectly well-defined dg coalgebra---it \emph{is} the Koszul dual coalgebra in the derived sense (Convention~\ref{conv:bar-coalgebra-identity})---but its cohomology does not concentrate in a single degree, so the classical Koszul dual algebra $\cA^!$ is not directly visible.  $I$-adic completion resolves this: the completed bar coalgebra $\widehat{\bar{B}}(\cA)$ recovers $\cA^!$ as a graded piece of a filtered inverse limit.
\end{observation}

\subsection{The completion solution}

\begin{idea}[\texorpdfstring{$I$}{I}-adic completion]\label{idea:I-adic}
The bar complex $\bar{B}(\mathcal{A})$ carries the $I$-adic (descending) filtration:
\[\bar{B}(\mathcal{A}) = I^0 \supseteq I^1 \supseteq I^2 \supseteq \cdots\]

where $I = \ker(\epsilon)$ is the coaugmentation coideal.

The $I$-adic completion is:
\[\widehat{\bar{B}}(\mathcal{A}) := \varprojlim_n \bar{B}(\mathcal{A})/I^n\]

This completed object is the candidate dual model in the completion
regime studied here.
\end{idea}

\section{The conilpotent filtration}

\subsection{Definition and properties}

\begin{definition}[Conilpotent filtration]\label{def:conilpotent-filtration}
For a coaugmented coalgebra $C$ with counit $\epsilon: C \to k$, coaugmentation $\eta: k \to C$, and coproduct $\Delta: C \to C \otimes C$, the \emph{coradical filtration} (or conilpotent filtration) is:
\begin{align*}
F_0C &= \operatorname{im}(\eta) = k \quad (\text{coaugmentation}) \\
F_nC &= \{c \in C : \Delta(c) \in F_{n-1}C \otimes C + C \otimes F_{n-1}C\} \quad (n \geq 1)
\end{align*}
Equivalently, using the reduced coproduct $\bar{\Delta}: \ker(\epsilon) \to \ker(\epsilon) \otimes \ker(\epsilon)$ and its iterates $\bar{\Delta}^{(n)}: \ker(\epsilon) \to \ker(\epsilon)^{\otimes(n+1)}$:
\[
F_nC = k \oplus \ker(\bar{\Delta}^{(n)}).
\]
In particular, $F_1C = k \oplus \ker(\bar{\Delta})$ consists of the scalars and the primitive elements.

This is an \emph{ascending} filtration: $F_0C \subset F_1C \subset F_2C \subset \cdots$.
A coalgebra is conilpotent if $\bigcup_n F_nC = C$, equivalently, if every element of $\ker(\epsilon)$ is annihilated by a sufficiently high iterate of the reduced coproduct: $\ker(\epsilon) = \bigcup_n \ker(\bar{\Delta}^{(n)})$.
\end{definition}

\begin{proposition}[Geometric manifestation; \ClaimStatusProvedHere]\label{prop:geom-conilpotent}
For the geometric bar complex:
\[\bar{B}^{\mathrm{geom}}_n(\mathcal{A}) = \Gamma(\overline{C}_{n+1}(X), \mathcal{A}^{\boxtimes(n+1)} \otimes \Omega^*_{\log})\]

The conilpotent filtration manifests as:
\begin{itemize}
\item $F_0$: Forms supported on the deepest stratum (all points colliding)
\item $F_1$: Forms supported on strata with at most one collision cluster
\item $F_n$: Forms supported on strata with at most $n$ nested collision levels
\end{itemize}

A form has collision depth~$n$ if it involves configurations where $n$ distinct groups of points simultaneously collide.
\end{proposition}

\begin{proof}
The Fulton--MacPherson compactification $\overline{C}_{n+1}(X)$ has a natural stratification indexed by trees~$T$ encoding collision patterns.  A stratum~$S_T$ corresponds to a configuration where the set $\{z_0, \ldots, z_n\}$ collides according to the tree~$T$: the leaves are the points, and each internal vertex represents a cluster of points that have collided.

Define the \emph{depth} of a tree $T$ as the maximal number of internal vertices on a root-to-leaf path.  The conilpotent filtration on $\bar{B}^{\mathrm{geom}}(\mathcal{A})$ is then:
\[
F_k \bar{B}^{\mathrm{geom}}_n = \bigoplus_{\substack{T \in \mathrm{Trees}(n+1) \\ \mathrm{depth}(T) \leq k}} \Gamma(S_T, \mathcal{A}^{\boxtimes(n+1)} \otimes \Omega^*_{\log}).
\]
This is ascending ($F_0 \subset F_1 \subset \cdots$) and exhaustive ($\bigcup_k F_k = \bar{B}^{\mathrm{geom}}_n$) because every tree has finite depth.

To verify this is the conilpotent filtration of Definition~\ref{def:conilpotent-filtration}: the reduced coproduct $\bar{\Delta}$ on $\bar{B}$ is the deconcatenation coproduct, which geometrically corresponds to ``splitting'' a configuration into two sub-configurations along a cut of the collision tree.  Cutting a tree at an internal vertex reduces the depth by at least one in each factor, so $\bar{\Delta}(F_k) \subset \sum_{i+j=k-1} F_i \otimes F_j$, matching Definition~\ref{def:conilpotent-filtration}.
\end{proof}

\subsection{Convergence of the completion}

\begin{theorem}[Completion convergence; \ClaimStatusProvedHere]\label{thm:completion-convergence}
For a chiral algebra $\mathcal{A}$ satisfying:
\begin{enumerate}
\item Finite generation over $\mathcal{D}_X$
\item Polynomial growth of structure constants
\item Finite-dimensional Hochschild cohomology: $\dim H^*(\mathcal{A}, \mathcal{A}) < \infty$
\end{enumerate}

The $I$-adic completion $\widehat{\bar{B}}(\mathcal{A})$ converges and defines a well-defined dg coalgebra.
\end{theorem}

\begin{proof}
Let $F^n \bar{B}(\mathcal{A})$ denote the \emph{descending} $I$-adic filtration (the sub-$\mathcal{D}_X$-module generated by elements of OPE order $\geq n$), so that $F^0 = \bar{B} \supset F^1 \supset F^2 \supset \cdots$.  (This is the $I$-adic filtration, not the ascending coradical filtration of Definition~\ref{def:conilpotent-filtration}.)  We verify that the inverse system $\{\bar{B}(\mathcal{A})/F^n\}_{n \geq 0}$ defines a convergent pro-coalgebra whose limit inherits a well-defined dg structure.

\emph{Step~1. Finite generation in each bidegree.}
Since $\mathcal{A}$ is finitely generated over~$\mathcal{D}_X$ by hypothesis, say by generators $a_1, \ldots, a_r$, the degree-$k$ component of the bar complex is:
\[
\bar{B}_k(\mathcal{A}) = \bigoplus_{i_1 + \cdots + i_{k+1} \leq N(k)} \mathcal{A}_{i_1} \otimes \cdots \otimes \mathcal{A}_{i_{k+1}} \otimes \Omega^k_{\log}(\overline{C}_{k+1}(X))
\]
where $N(k)$ is a bound coming from the conformal weight grading.  The filtration level $F^n \cap \bar{B}_k$ is a sub-$\mathcal{D}_X$-module of a finitely generated module, hence finitely generated (since $\mathcal{D}_X$ is Noetherian on a smooth algebraic curve~$X$).

\emph{Step~2. Polynomial growth controls the associated graded.}
The polynomial growth hypothesis states that the OPE coefficients $C^c_{ab}(n)$ (defined by $a_{(n)}b = \sum_c C^c_{ab}(n)\, c$) satisfy $|C^c_{ab}(n)| \leq P(n)$ for a polynomial~$P$.  This implies that the bar differential $d_{\bar{B}}$, which involves residues of OPEs, preserves the filtration ($d_{\bar{B}}(F^n) \subset F^n$) with the induced map on graded pieces $\mathrm{gr}^n \bar{B} \to \mathrm{gr}^n \bar{B}$ bounded polynomially in~$n$.  Therefore:
\[
\dim\bigl((F^n/F^{n+1}) \cap \bar{B}_k\bigr) \leq P(n) \cdot Q(k)
\]
for polynomials $P, Q$ depending only on the generating data of~$\mathcal{A}$.

\emph{Step~3. Finiteness ensures differential compatibility.}
The finiteness condition $\dim H^*(\mathcal{A}, \mathcal{A}) < \infty$ guarantees that the obstruction theory for extending the differential from $\bar{B}/F_n$ to $\bar{B}/F_{n+1}$ is unobstructed.  Specifically, the obstruction to lifting $d_{\bar{B}} \bmod F^n$ to $d_{\bar{B}} \bmod F^{n+1}$ lives in $H^2(\mathcal{A}, F^n/F^{n+1})$, which is finite-dimensional and eventually zero by the polynomial growth bound.  (This step uses the polynomial growth hypothesis of~(H1) to ensure eventual stabilization; see Remark~\ref{rem:completion-convergence-frontier} for the proof-density discussion.)

\emph{Step~4. Mittag-Leffler condition and inverse limit.}
The transition maps $\bar{B}_k/F^{n+1} \twoheadrightarrow \bar{B}_k/F^n$ are surjections (since $F^{n+1} \subset F^n$).  For any inverse system of surjections, the Mittag-Leffler condition is automatically satisfied, which implies $\varprojlim^1_n (\bar{B}_k/F^n) = 0$.  (The polynomial growth bound on $\dim(F^n/F^{n+1})_k$ ensures that each quotient is finite-dimensional, so the inverse system is in the category of finite-dimensional vector spaces where the inverse limit is exact.)  Therefore the inverse limit
\[
\widehat{\bar{B}}_k(\mathcal{A}) = \varprojlim_n \bar{B}_k(\mathcal{A})/F^n
\]
exists and is exact.  The dg coalgebra structure passes to the limit because $d_{\bar{B}}$ and $\Delta$ are continuous for the $F^\bullet$-adic topology.
\end{proof}

\begin{remark}[Frontier status]\label{rem:completion-convergence-frontier}
Step~3 above compresses a non-trivial claim: that the obstruction
groups $H^2(\mathcal{A}, F^n/F^{n+1})$ are eventually zero under the
polynomial growth hypothesis.  The argument invokes the finite
Hochschild cohomology assumption to control the total obstruction space,
and the polynomial growth bound to show each graded piece $F^n/F^{n+1}$
becomes negligible for large~$n$.  A fully expanded proof would require
an explicit filtration-degree estimate showing that the polynomial bound
on OPE coefficients forces $H^2(\mathcal{A}, F^n/F^{n+1}) = 0$ for
$n \gg 0$.  This places the result in the completion frontier regime of
Remark~\ref{rem:existence-regimes}(ii), not in the strict finite-type
core.
\end{remark}

\section{Koszul duality via completion}

\subsection{The main construction}

\begin{construction}[Completed Koszul dual]\label{const:completed-koszul}
For non-quadratic chiral algebras, the intrinsic
Definition~\ref{def:koszul-dual-quadratic} does not apply.
Define the \emph{completed Koszul dual coalgebra} as:
\[\mathcal{A}^! := \widehat{\bar{B}}(\mathcal{A})\]
the $I$-adic completion of the geometric bar complex.
This is a completed dg \emph{coalgebra}; the Koszul dual algebra
(if it exists) is obtained by continuous linear dualizing under
topological finiteness hypotheses.

The cobar construction then gives:
\[\Omega(\mathcal{A}^!) := (T(s^{-1}\overline{\mathcal{A}^!}), d_\Omega)\]
where $T$ denotes the tensor algebra, $s^{-1}$ the desuspension, $\overline{\mathcal{A}^!}$ the coaugmentation coideal, and $d_\Omega$ is the differential induced by the coproduct of $\mathcal{A}^!$.
\end{construction}

\begin{theorem}[Completed bar-cobar duality; \ClaimStatusProvedHere]\label{thm:completed-bar-cobar}
For a chiral algebra $\mathcal{A}$ satisfying the convergence criteria of
Theorem~\textup{\ref{thm:completion-convergence}}, there is a quasi-isomorphism:
\[\Omega(\widehat{\bar{B}}(\mathcal{A})) \simeq \mathcal{A}\]

This is the completion-regime analogue of the finite-type bar-cobar
inversion of the core monograph (Main Theorem~B), valid under the
finite-generation and polynomial-growth hypotheses stated above.
See Remark~\textup{\ref{rem:completed-bar-cobar-frontier}} for the
frontier status of the argument.
\end{theorem}

\begin{proof}
The proof extends the uncompleted bar-cobar inversion (Main Theorem~B) via a comparison argument using the completion filtration.

\emph{Step~1. Filtration compatibility.}
The descending bar-degree filtration $\{F^n\bar{B}(\mathcal{A})\}_{n \geq 0}$
(where $F^n = \bigoplus_{k \geq n} \bar{B}^{\mathrm{ch}}_k(\mathcal{A})$)
induces a completion:
$\widehat{\bar{B}}(\mathcal{A}) = \varprojlim_n \bar{B}(\mathcal{A})/F^n$.
The bar differential $d_{\bar{B}}$ is compatible with this filtration
(it increases bar degree by at most~1),
so $\widehat{\bar{B}}(\mathcal{A})$ is a well-defined dg coalgebra.

\emph{Step~2. Cobar functor on the completion.}
The cobar functor $\Omega$ applied to $\widehat{\bar{B}}(\mathcal{A})$ produces a dg algebra whose underlying graded space is:
\[
\Omega(\widehat{\bar{B}}(\mathcal{A})) = T(s^{-1}\overline{\widehat{\bar{B}}(\mathcal{A})})
\]
with differential $d_\Omega$ encoding the coproduct of $\widehat{\bar{B}}(\mathcal{A})$.  We must show this is quasi-isomorphic to~$\mathcal{A}$.

\emph{Step~3. Comparison with the uncompleted version.}
The canonical map $\bar{B}(\mathcal{A}) \to \widehat{\bar{B}}(\mathcal{A})$ induces a map of cobar complexes:
\[
\Omega(\bar{B}(\mathcal{A})) \longrightarrow \Omega(\widehat{\bar{B}}(\mathcal{A}))
\]
induced by the canonical inclusion $\bar{B}(\mathcal{A}) \hookrightarrow \widehat{\bar{B}}(\mathcal{A})$ (the cobar functor is covariant).  We show this map is a quasi-isomorphism using Step~4.  By the uncompleted bar-cobar inversion theorem, $\Omega(\bar{B}(\mathcal{A})) \xrightarrow{\sim} \mathcal{A}$ is a quasi-isomorphism.

\emph{Step~4. Completion does not change homotopy type.}
The convergence hypotheses (Theorem~\ref{thm:completion-convergence}) ensure that the inverse system $\{\bar{B}(\mathcal{A})/F^n\}_{n \geq 0}$ (with $F^n$ the descending $I$-adic filtration) satisfies the Mittag-Leffler condition.  Indeed, the polynomial growth of structure constants implies that for each bar complex degree~$k$, the quotients $(F^n \bar{B}_k)/(F^{n+1} \bar{B}_k)$ stabilize for $n$ sufficiently large relative to~$k$.  Therefore $\varprojlim^1 = 0$, and the natural map
\[
H^*(\widehat{\bar{B}}(\mathcal{A})) \xrightarrow{\;\cong\;} \varprojlim_n H^*(\bar{B}(\mathcal{A})/F^n)
\]
is an isomorphism.  Since each finite-level quotient $\bar{B}(\mathcal{A})/F_n$ is a conilpotent coalgebra (by Positselski \cite{Positselski11}, Theorem~3.5), the cobar construction on it recovers~$\mathcal{A}/I^n$ (where $I$ is the augmentation ideal of~$\mathcal{A}$), and passing to the limit gives $\Omega(\widehat{\bar{B}}(\mathcal{A})) \simeq \varprojlim_n (\mathcal{A}/I^n) = \mathcal{A}$, the last equality holding because $\mathcal{A}$ is $I$-adically complete under the convergence hypotheses.
\end{proof}

\begin{remark}[Frontier status]\label{rem:completed-bar-cobar-frontier}
Step~4 above uses finite-stage bar-cobar inversion: the claim that
$\Omega(\bar{B}(\mathcal{A})/F_n) \simeq \mathcal{A}/I^n$ for each~$n$
requires that each truncation $\mathcal{A}/I^n$ satisfies the hypotheses
of the uncompleted bar-cobar inversion theorem (Theorem~B).  For the
standard Master Table families, these finite-stage quotients are
finite-dimensional and hence automatically on the Koszul locus.  For
algebras outside these families, this finite-stage compatibility is an
additional hypothesis.  This places the result in the completion
frontier regime of Remark~\ref{rem:existence-regimes}(ii).
\end{remark}

\subsection{Essential image}

\begin{theorem}[Characterization of Koszul duals; \ClaimStatusProvedHere]\label{thm:koszul-dual-characterization}
A completed coalgebra $\widehat{\mathcal{C}}$ arises as $\widehat{\mathcal{C}} \simeq \widehat{\bar{B}}(\mathcal{A})$ for some chiral algebra $\mathcal{A}$ satisfying the convergence criteria of Theorem~\textup{\ref{thm:completion-convergence}} if and only if:
\begin{enumerate}
\item $\widehat{\mathcal{C}}$ is a completed conilpotent coalgebra
\item The coderivation $d$ is compatible with the completion topology
\item There exist ``generating cogenerators'' $V^* \hookrightarrow F_1\widehat{\mathcal{C}}/F_0\widehat{\mathcal{C}}$ compatible with collision data on $\overline{C}_n(X)$
\item The spectral sequence $E_1 = \operatorname{gr}(\widehat{\mathcal{C}})$ degenerates at $E_2$
\end{enumerate}
\end{theorem}

\begin{proof}
\emph{Forward direction} ($\widehat{\mathcal{C}} = \widehat{\bar{B}}(\mathcal{A}) \Rightarrow$ (1)--(4)):

\emph{Condition (1).} The bar complex $\bar{B}(\mathcal{A})$ is conilpotent by construction: the deconcatenation coproduct $\bar{\Delta}$ on $\bar{B}_n$ reduces bar degree, so $\bar{\Delta}^{(n)}(\bar{B}_n) = 0$.  Proposition~\ref{prop:geom-conilpotent} gives the geometric manifestation.  The completion
$\widehat{\bar{B}}(\mathcal{A}) = \varprojlim_n \bar{B}(\mathcal{A})/F_n$
inherits conilpotency in the completed sense: every element of $\ker(\hat{\epsilon})$ lies in some~$F_n$.

\emph{Condition (2).} The bar differential $d_{\bar{B}}$ is defined via residues of OPEs, which involve OPE coefficients at fixed singularity order.  Since $d_{\bar{B}}$ maps $F_n$ to $F_{n-1}$ (collision depth decreases by at most one under residue extraction), it is continuous for the $F$-adic topology.  This was verified in Step~1 of the proof of Theorem~\ref{thm:completed-bar-cobar}.

\emph{Condition (3).} The generators $a_1, \ldots, a_r$ of $\mathcal{A}$ embed as cogenerators of $\bar{B}(\mathcal{A})$ via the canonical inclusion $\mathcal{A} \hookrightarrow \bar{B}^1(\mathcal{A})$ (the bar-degree-one component).  These elements lie in $F_1/F_0$ and satisfy collision compatibility: the coproduct $\bar{\Delta}$ on a cogenerator $a_i \in \bar{B}^1$ vanishes (they are primitive), while the differential $d(a_i \otimes a_j \otimes \eta_{12}) = \operatorname{Res}_{z_1 = z_2}[a_i(z_1)a_j(z_2)]$ recovers the OPE data.  This is precisely the compatibility with collision divisors on $\overline{C}_n(X)$.

\emph{Condition (4).} The spectral sequence associated
to the conilpotent filtration has
$E_1^{p,q} = H^q(\operatorname{gr}^p
\widehat{\bar{B}}(\mathcal{A}))$.
The $E_1$ page computes the cohomology of the associated graded
coalgebra, which for the bar complex of $\mathcal{A}$ is determined by
the quadratic part of the relations.  The $E_2$ degeneration holds
because the convergence criteria of
Theorem~\ref{thm:completion-convergence} (finite generation +
polynomial growth) ensure that $E_2^{p,q}$ is finite-dimensional in
each bidegree, and for degree reasons ($E_r$ differentials have bidegree
$(r, 1-r)$) all higher differentials $d_r$ for $r \geq 3$ vanish on a
range that exhausts the complex as $n \to \infty$.

\medskip
\emph{Backward direction} ((1)--(4) $\Rightarrow$ $\widehat{\mathcal{C}} \simeq \widehat{\bar{B}}(\mathcal{A})$):

Given $\widehat{\mathcal{C}}$ satisfying (1)--(4), define $\mathcal{A} := \Omega(\widehat{\mathcal{C}})$, the cobar construction on the completed coalgebra.

\emph{Step 1: Well-definedness.}  Conditions (1) and (2) ensure that $\Omega(\widehat{\mathcal{C}}) = (T(s^{-1}\overline{\widehat{\mathcal{C}}}), d_\Omega)$ is a well-defined dg algebra: the conilpotency guarantees that the cobar differential $d_\Omega$ (induced by $\bar{\Delta}$ and $d$) involves only finite sums, and continuity of $d$ ensures $d_\Omega^2 = 0$ on the completed tensor algebra.

\emph{Step 2: Chiral structure.}  Condition (3) provides generators $V^* \subset F_1/F_0$ whose collision compatibility endows $\mathcal{A} = \Omega(\widehat{\mathcal{C}})$ with a chiral algebra structure: the cobar differential encodes OPE data via $d_\Omega(s^{-1}v)(z_1, z_2) = \operatorname{Res}_{z_1 = z_2}[\cdots]$, and the associativity of the chiral product follows from $d_\Omega^2 = 0$.  The collision compatibility ensures that these OPEs satisfy the locality axiom (singularities only along diagonals).

\emph{Step 3: Recovery.}  By Theorem~\ref{thm:completed-bar-cobar}, $\Omega(\widehat{\bar{B}}(\mathcal{A})) \simeq \mathcal{A}$.  We must show $\widehat{\bar{B}}(\mathcal{A}) \simeq \widehat{\mathcal{C}}$.  The bar-cobar adjunction provides a unit $\eta: \widehat{\mathcal{C}} \to \widehat{\bar{B}}(\Omega(\widehat{\mathcal{C}}))$.  Condition (4) ensures this is a quasi-isomorphism: the $E_2$ degeneration means the bar spectral sequence for $\mathcal{A} = \Omega(\widehat{\mathcal{C}})$ collapses, so $\widehat{\bar{B}}(\mathcal{A})$ has the same cohomology as $\widehat{\mathcal{C}}$.  Since both are complete and conilpotent, a cohomology isomorphism between completed conilpotent coalgebras is a quasi-isomorphism of dg coalgebras (by the comparison theorem for complete filtered complexes, cf.\ Positselski \cite{Positselski11}).
\end{proof}

\begin{remark}[Frontier status]\label{rem:koszul-dual-characterization-frontier}
The backward direction (Step~3) invokes
Theorem~\ref{thm:completed-bar-cobar}, whose frontier status is
recorded in Remark~\ref{rem:completed-bar-cobar-frontier}.  The
essential-image characterization therefore inherits the same
frontier standing: the result is valid under the convergence
hypotheses of Theorem~\ref{thm:completion-convergence}, but those
hypotheses are part of the completion frontier regime, not the strict
finite-type core.
\end{remark}

\section{Connection to Beilinson--Drinfeld}

\subsection{Filtered chiral algebras}

\begin{remark}[BD framework]\label{rem:BD-filtered}
Beilinson--Drinfeld (BD 4.2) develop the theory of filtered chiral algebras with I-adic topology. Their framework provides:

\begin{itemize}
\item Filtration by augmentation ideal I
\item Completion: $\widehat{\mathcal{A}} = \varprojlim_n \mathcal{A}/I^n$
\item Associated graded: $\text{gr}(\mathcal{A}) = \bigoplus_n I^n/I^{n+1}$
\end{itemize}

Our nilpotent completion of the bar complex is an instance of this framework:
\[\widehat{\bar{B}}(\mathcal{A}) = \varprojlim_n \bar{B}(\mathcal{A})/I^n\]
is the BD-style completion applied to the bar complex.
\end{remark}

\subsection{Chiral homology completion}

\begin{theorem}[BD chiral homology \cite{BD04}; \ClaimStatusProvedElsewhere]\label{thm:BD-chiral-homology}
\textup{(Beilinson--Drinfeld 4.7)}

For a chiral algebra $\mathcal{A}$ and curve X, BD define completed chiral homology:
\[\widehat{H}_*^{\text{ch}}(X, \mathcal{A}) := \varprojlim_n H_*^{\text{ch}}(X, \mathcal{A}/I^n)\]

This agrees with our completed bar construction:
\[\widehat{H}_*^{\text{ch}}(X, \mathcal{A}) \cong H_*(\widehat{\bar{B}}(\mathcal{A}))\]
\end{theorem}

\section{Physical interpretation: renormalization}

\subsection{Perturbative expansion}

\begin{interpretation}[Completion as perturbative expansion]\label{interp:completion-renorm}
In the Poisson sigma model on a surface~$\Sigma_g$, the filtration $F^n$ of the bar complex corresponds to the loop expansion truncated at~$n$ loops: the finite truncation $\bar{B}(\mathcal{A})/I^n$ captures contributions from Feynman diagrams with at most~$n$ internal edges, and the completion $\widehat{\bar{B}}(\mathcal{A}) = \varprojlim_n \bar{B}(\mathcal{A})/I^n$ is the all-orders perturbative expansion.  The passage from finite truncation to inverse limit is the algebraic counterpart of removing the UV cutoff, and the convergence criteria of Theorem~\ref{thm:completion-convergence} ensure the limit is well-defined.
\end{interpretation}

\subsection{Feynman diagram expansion}

\begin{remark}[Feynman diagrams and conilpotent degree]\label{rem:feynman-conilpotent}
The conilpotent filtration levels $F^n/F^{n+1}$ correspond, under the Feynman diagram identification of Chapter~\ref{ch:feynman}, to graphs with exactly~$n$ internal edges.  In particular: $F^0/F^1$ consists of vacuum contributions (no internal edges), $F^1/F^2$ of tree-level propagators (single internal edge), and $F^n/F^{n+1}$ of diagrams with~$n$ internal edges.  The symmetry factors arising in the perturbative expansion are the orders of the automorphism groups of the corresponding graphs.  The completion $\widehat{\bar{B}}(\mathcal{A}) = \varprojlim_n \bar{B}(\mathcal{A})/I^n$ amounts to summing over all graph topologies---a formal power series that requires the $I$-adic topology to converge.
\end{remark}

\section{Explicit examples}

\subsection{Example 1: Virasoro algebra}

\begin{example}[Virasoro completion]\label{ex:virasoro-completion}
The Virasoro algebra has generators $L_n$ with OPE:
\[T(z)T(w) \sim \frac{c/2}{(z-w)^4} + \frac{2T(w)}{(z-w)^2} + \frac{\partial T(w)}{z-w}\]

The quartic pole means the naive bar complex has:
\[\bar{B}_2(\text{Vir}) \ni L_n \otimes L_m \otimes (z-w)^{-4} dz \wedge dw\]

This requires completion to handle all powers of $(z-w)^{-k}$ for $k \geq 4$.

\emph{Filtration.}
\begin{align*}
F_0 &= \text{All L-generators} \\
F_1 &= \text{Composite fields like } :LL:, :LLL:, \ldots \\
F_2 &= \text{Double composites} \\
&\vdots
\end{align*}

\emph{Completion.}
\[\widehat{\bar{B}}(\text{Vir}) = \varprojlim_n \bar{B}(\text{Vir})/I^n\]

Includes all formal power series in composite operators.

\emph{Koszul dual.}
\[\text{Vir}^! = \widehat{\bar{B}}(\text{Vir})\]
is a completed coalgebra with infinitely many cogenerators corresponding to the composite fields.
\end{example}

\subsection{\texorpdfstring{Example 2: $W_3$ algebra}{Example 2: W-3 algebra}}

\begin{example}[\texorpdfstring{$W_3$}{W3} completion]\label{ex:w3-completion}
The W₃ algebra has generators T (spin 2) and W (spin 3) with OPE:
\[W(z)W(w) \sim \frac{c/3}{(z-w)^6} + \cdots\]

The sixth-order pole requires even more careful completion.

\emph{Convergence.}
The completion converges because:
\begin{enumerate}
\item Finite generation: Two generators T, W
\item Polynomial growth: Structure constants grow like $n^k$ for fixed k
\item Formal smoothness: $\dim H^*(\mathcal{W}_3, \mathcal{W}_3) < \infty$
\end{enumerate}

\emph{Completed bar complex.}
\[\widehat{\bar{B}}(\mathcal{W}_3) = \varprojlim_{n} \bar{B}(\mathcal{W}_3)/I^n\]

includes all formal composite fields built from T and W.

\emph{Geometric interpretation.}
At the geometric level, completion corresponds to summing over all collision patterns on configuration spaces, including arbitrarily nested collisions.
\end{example}

\subsection{Example 3: Heisenberg (quadratic case for comparison)}

\begin{example}[Heisenberg: no completion needed]\label{ex:heisenberg-no-completion}
The Heisenberg algebra is quadratic:
\[a(z)a(w) \sim \frac{k}{(z-w)^2}\]

The bar complex has:
\[\bar{B}_n(\mathcal{H}_k) = \mathcal{H}_k^{\otimes n} \otimes \Omega^*_{\log}(\overline{C}_{n+1}(X))\]

with finitely many terms in each degree.

\emph{No completion required.}
\[\bar{B}(\mathcal{H}_k) = \widehat{\bar{B}}(\mathcal{H}_k)\]

The Heisenberg algebra is Koszul (its quadratic part $\operatorname{Sym}(V)$ is
Koszul with dual $\Lambda(V^*)$, a finite-dimensional coalgebra).
Since the Koszul dual coalgebra is finite-dimensional, it is automatically
conilpotent, and the bar complex is quasi-isomorphic to this finite model
without any completion.  This is the general mechanism: for Koszul algebras, the
bar complex admits a small, conilpotent model, so completion is unnecessary.
\end{example}

\section{Computational methods}

\subsection{Computing the completion}

\begin{remark}[Computing the truncation \texorpdfstring{$\widehat{\bar{B}}(\mathcal{A})/I^N$}{B-hat(A)/I\textasciicircum N}]\label{alg:compute-completion}
Given a chiral algebra $\mathcal{A}$ and truncation level $N$, the quotient $\bar{B}(\mathcal{A})/I^N$ is computed as follows.
\begin{enumerate}
\item \emph{Generate basis.} Degree 0: vacuum; degree 1: generators of $\mathcal{A}$; degree $k$: $k$-fold tensor products.
\item \emph{Compute differential.} Use the OPE to compute $d(a_1 \otimes \cdots \otimes a_k)$ by extracting residues at collisions and summing over all collision patterns.
\item \emph{Filter by conilpotent degree.} Compute the coproduct iteratively, determine which elements lie in $I^n$ for each $n$, and truncate at level $N$.
\end{enumerate}
\end{remark}

\subsection{Convergence criteria in practice}

\begin{proposition}[Convergence of the completion \cite{BD04, FG12}; \ClaimStatusProvedElsewhere]\label{prop:practical-convergence}
For a chiral algebra $\mathcal{A}$ graded by conformal weight, the completion
$\widehat{\bar{B}}(\mathcal{A})$ is well-defined (the projective limit stabilizes
in each total weight) whenever:
\begin{enumerate}
\item $\mathcal{A}$ is \emph{finitely generated}; the bar complex at bar degree $n$
has total weight $\geq n \cdot h_{\min}$ where $h_{\min} = \min_i h_i > 0$,
so each weight space is automatically finite-dimensional.  No completion is needed
beyond finite truncation.
\item $\mathcal{A}$ has \emph{countably many generators} of weights
$h_1 \leq h_2 \leq \cdots \to \infty$; the completion converges if the weights grow
fast enough that each total-weight subspace of the bar complex is finite-dimensional.
\end{enumerate}

\emph{Examples.}
\begin{itemize}
\item Virasoro ($h = 2$), $W_3$ ($h = 2, 3$): finitely generated, completion automatic.
\item $\mathcal{W}_{1+\infty}$: generators in all integer weights $h = 1, 2, 3, \ldots$;
completion requires care but converges since $h_n \to \infty$.
\end{itemize}
\end{proposition}

\section{Connection to Costello--Gwilliam}

\subsection{Renormalization and factorization}

\begin{theorem}[Costello--Gwilliam renormalization \cite{CG17}; \ClaimStatusProvedElsewhere]\label{thm:CG-renorm}
\textup{(Costello--Gwilliam, Volume 2)}

For a perturbative quantum field theory, the effective interaction satisfies:
\[I_{\text{eff}} = I_{\text{classical}} + \sum_{n=1}^{\infty} \hbar^n I^{(n)}\]

where $I^{(n)}$ are n-loop quantum corrections.

This is the structure of our $I$-adic completion:
\[\widehat{\bar{B}}(\mathcal{A}) = \varprojlim_{n} \bar{B}(\mathcal{A})/I^n\]

\end{theorem}

\begin{remark}[Physical interpretation]
Under the Feynman diagram identification of Chapter~\ref{ch:feynman}, the $I$-adic completion topology is the algebraic counterpart of the loop expansion: convergence of the inverse system $\{\bar{B}(\mathcal{A})/I^n\}$ ensures that the formal sum over all loop orders defines a well-defined dg coalgebra.  The analogy with renormalization group flow is suggestive---both organize quantum corrections by a filtration parameter---but the precise relationship requires the graph complex technology developed in that chapter.
\end{remark}

\section{Higher genus corrections}

\subsection{Genus expansion of completion}

\begin{remark}[Genus-graded completion]\label{rem:genus-completion}
At higher genus, the completion becomes even more essential. The bar complex receives contributions from all genera:
\[\widehat{\bar{B}}_g(\mathcal{A}) = \varprojlim_n \bar{B}_g(\mathcal{A})/I^n\]

The genus-g contribution involves:
\begin{itemize}
\item Configuration spaces on Riemann surfaces of genus g
\item Modular forms for genus g
\item Period integrals over $\mathcal{M}_g$
\end{itemize}

The completion sums over all genera as a formal Laurent series in $\hbar$:
\[\widehat{\bar{B}}(\mathcal{A}) = \sum_{g=0}^{\infty} \hbar^{2g-2} \widehat{\bar{B}}_g(\mathcal{A}) = \hbar^{-2}\widehat{\bar{B}}_0(\mathcal{A}) + \widehat{\bar{B}}_1(\mathcal{A}) + \hbar^2 \widehat{\bar{B}}_2(\mathcal{A}) + \cdots\]

This is the string theory expansion (the $\hbar^{-2}$ corresponds to the sphere/tree-level contribution).
\end{remark}

\section{Summary}

\begin{remark}[Completed frontier package]\label{rem:nilpotent-main}
This appendix should be read as a frontier package, not as a blanket
extension of the proved core.  What is currently in hand is:
\begin{enumerate}
\item the finite-stage convergence and geometry of the completed bar
      tower
      \textup{(}Theorem~\ref{thm:completion-convergence},
      Proposition~\ref{prop:geom-conilpotent}\textup{)};
\item the inverse-limit criteria used by the standard
      $\mathcal{W}_\infty$ and Yangian towers in the MC4 package of
      Chapter~\ref{chap:concordance};
\item the compatibility of these completed models with the
      Beilinson--Drinfeld filtered framework
      \textup{(}Theorem~\ref{thm:BD-chiral-homology}\textup{)}.
\end{enumerate}
What is \emph{not} yet elevated to monograph-wide theorematic status is
the universal claim that finite generation plus polynomial growth alone
suffice for every non-quadratic chiral algebra.  The remaining missing
steps are exactly the global obstruction-vanishing argument, the
completed $E_2$-degeneration argument, and the finite-stage
compatibility needed to pass from the standard towers to a universal
inverse-limit theorem.
\end{remark}
